\begin{center}
    {\large \textbf{Epilogo}}
\end{center}

\noindent \newline \newline


È trascorsa una settimana dal nostro burrascoso ritorno ed è giunto il momento di tirare le somme.

Il mattino dopo essere arrivati a casa, dopo aver tentato invano di riposarmi, ho ricominciato subito a lavorare — per fortuna non in presenza — facendo molta fatica a concentrarmi, sia per la stanchezza fisica, sia perché la testa non riusciva a star ferma e continuava a viaggiare. Come al ritorno dal Messico, quando seduto davanti al computer in ufficio la mia testa era ancora nei cieli sopra Teotihuacan, anche qui mi rivedevo con Oshi a camminare tra le risaie, in kayak a sfidare vento e onde lungo le coste della Gili Asahan, a immergermi con le tartarughe davanti alla spiaggia di Gili Trawangan, a scappare dalle scimmie che volevano il mio sacchetto di patatine nella Monkey Forest di Ubud, o a chiacchierare con Nyoman su come Bali sia cambiata nel corso degli anni.

Così come per il Messico, sento di poter dire che non abbiamo solo visitato l'Indonesia: l'abbiamo vissuta. Abbiamo cercato il più possibile esperienze vere, dove la parola "turismo" ancora fatica ad arrivare, nel bene e nel male.

\begin{center}
    % Decorative line
    \noindent\rule{0.6\textwidth}{0.4pt}
\end{center}

È così che ci siamo accorti di come, laddove le città sono più ricche, più commerciali, più mete consolidate del turismo occidentale, la popolazione locale sorrida di meno.

Fin dal primo giorno nel Tetebatu — sicuramente la regione più povera che abbiamo visitato — i sorrisi dei locali ci hanno tenuto compagnia e contagiato, così come il loro umore sempre gioviale e mai altalenante. Sulla Gili Asahan, una meta sconosciuta ai più, un'isola abitata da ottanta persone, non hanno smesso di sorriderci nemmeno quando ce ne siamo andati, continuando a sbracciarsi per salutare anche quando la nostra barca era ormai lontana dalla costa. La "Gili" che non appartiene a "Le Gili" — gili in indonesiano significa isola, ma quando si parla delle Gili tutti si riferiscono alle tre isole a nord-ovest di Lombok, spesso meta di qualche giorno di mare per chi associa l'Indonesia alla sola Bali. Mi spiace dirlo, ma se sei stato solo sulle Gili, non puoi dire di essere stato a Lombok.

Già a Kuta Lombok, meta preferita dei surfisti australiani, avevamo notato un mutamento nelle persone, come se un'ombra fosse calata sui loro volti. Qui abbiamo iniziato a guardarci dagli scam delle guide turistiche — dov'è il mio pugnale di Lombok? — dai sapori e dai prezzi occidentali del Novotel — vogliamo parlare della carbonara con panna e prosciutto? — riuscendo però ancora a vivere la cultura locale (il Betel Quid sono certo che Elisa non lo dimenticherà tanto facilmente) e a sentirci vicini ai suoi abitanti. Sento ancora Rudy per le nostre lezioni di italiano, e di recente mi ha mandato questo messaggio: "Later when you have friends want to have holiday in Lombok, let me know my teacher. They can stay in my house so they don't have to pay much for their hotel. My house only ten minutes from the airport."

La parte di vacanza sulla Gili Trawangan ha interrotto tale discesa, grazie a una serie di fattori determinanti che ci hanno permesso di vivere l'isola come abbiamo imparato a fare dagli indonesiani conosciuti a Lombok. L'aver incontrato Enrico e Cristina, i giochi per la festa dell'indipendenza indonesiana, l'atmosfera da paradiso tropicale e le nuotate con le tartarughe hanno sicuramente contribuito ad allontanarci dal vero mood della Gili T — la Party Island delle feste in discoteca, dei Pub Crawl e dei "mushrooms, marijuana, cocaine?" sussurrati lungo la strada.

Da quando ci siamo trasferiti a Bali, abbiamo notato un cambiamento nei sorrisi della gente del posto: prima erano semplici bocche serrate, espressioni di chi lavora per abitudine più che per piacere, poi sono diventati visi imbronciati, di persone presenti fisicamente ma con la mente altrove. Ubud, seppur ricca di storia e tradizione, è ormai una città occidentalizzata: punto nevralgico al centro della rete dei templi indù, ma per trovare una libreria che non venda solo libri internazionali bisogna allontanarsi di almeno un'ora di macchina.

Tuttavia, il passato è ancora presente, se lo si sa cercare. Si può gustare una cena nella cornice tradizionale del Caffè Wayan — che, a essere onesti, risale solo agli anni '80, ma almeno sembra aver resistito all'onda di occidentalizzazione dell'ultimo decennio. Oppure si può attraversare la città cercando di proteggere un pacchetto di patatine dalle scimmie di Monkey Forest, che ancora si sentono padrone di Ubud. A pochi chilometri di distanza è possibile visitare antichi templi dall'atmosfera unica — parlo di Samuan Tiga, dove il timido volontario ha strappato solo sei biglietti in un giorno, e non di Goa Gajah, dove devi scattare la foto della Grotta dell'Elefante in fretta, prima che un turista si metta in posa davanti.

In tre giorni siamo riusciti a visitare sette templi a Bali e, se dovessimo fare una classifica, sono certo che Elisa sarebbe d'accordo con me: i templi meno turistici sono i migliori, ad eccezione forse del meraviglioso Tirta Gangga, con le sue immense piscine e le alte fontane a più livelli. L'essere stati i primi a entrare ha aiutato a farcelo apprezzare maggiormente, poiché già alle nove del mattino il luogo cominciava a riempirsi di turisti intenti a fotografarsi mentre sovralimentavano i pesci fino a farli quasi scoppiare.

Kuta, infine, è stato il punto culminante di questo percorso verso l'occidentalizzazione: fast-food, centri commerciali, piscine con musica house, armi giocattolo, centri balneari discutibili, e soprattutto il monito di Nyoman: "Kuta is not Ubud." Tutti questi fattori, nel complesso, ci hanno fatto sentire un po' meno il dispiacere di dover tornare a casa.

\begin{center}
    % Decorative line
    \noindent\rule{0.6\textwidth}{0.4pt}
\end{center}

La Bali di oggi non è la Bali di dieci anni fa. Lombok forse si salva ancora — presumo che già nel 2010 le Gili fossero meta turistica per musica e feste, spesso viste come estensione della vacanza a Bali — ma abbiamo capito che il turismo trash sta piano piano arrivando anche qui. Anche i nativi lo hanno capito e sembra che l'ago della bilancia oscilli tra il marciarci sopra e il ripudiare i turisti. Negli anni i villaggi Sasak si stanno spopolando sempre di più, rimanendo infine ricostruzioni di quello che un tempo erano le loro vere abitazioni: musei a cielo aperto per turisti accompagnati da guide felici di chiedere mance sempre più alte, e allo stesso tempo risentiti per essere stati allontanati dalle loro case e dalle loro tradizioni per inseguire il progresso.

Forse tutto questo non succederà dall'oggi al domani — più probabilmente sarà un processo lento e naturale, che porterà anche benefici e cultura. Non è detto che solo "noi abbiamo tanto da imparare da loro", come ha spesso affermato Elisa durante questo viaggio. È possibile, invece, che anche loro possano trarre beneficio dalla nostra influenza, non solo in termini economici. Dal punto di vista sanitario, nella prevenzione, nella conservazione degli alimenti e persino nel modo in cui trattiamo gli animali, possiamo offrire un contributo significativo.

Quest'ultimo aspetto meriterebbe un approfondimento a parte. Durante queste due settimane abbiamo osservato come gli animali vengano considerati esclusivamente come mezzi di sostentamento, ricevendo poca cura e attenzione. I cavalli che trainavano i calessi sulle isole Gili erano costantemente legati ai loro carretti, anche di notte, sempre in movimento e con un sacchetto attaccato per raccogliere gli escrementi, costretti a portare in giro i propri bisogni per tutto il giorno. I cani, spesso randagi e raramente con collare e guinzaglio, erano lasciati a grattarsi fino a ferirsi, arrivando a esporre la carne viva, ignorati o scacciati dai locali come fossero dei lebbrosi. E poi ci sono gli animali quasi torturati a fini di lucro: le carpe sovralimentate a Tirta Gangga per qualche foto su Instagram, gli zibetti costretti a mangiare solo chicchi di caffè per produrre il costoso caffè Luwak, o i poveri pulcini colorati del mercato di Sukawati, che ancora ricordo con gli occhi chiusi e il corpo tremante.

\begin{center}
    % Decorative line
    \noindent\rule{0.6\textwidth}{0.4pt}
\end{center}

Un'altra riflessione riguarda la lingua. Come in Messico, dove la conoscenza anche minima dello spagnolo ci ha permesso di avvicinarci ai locali, qui conoscere l'inglese — soprattutto nella sua forma più informale — è stato qui un enorme vantaggio. Mi sono reso conto di quanto l'inglese colloquiale faciliti la costruzione di un rapporto più autentico con le persone in ogni situazione.

Per esempio, quando siamo andati sulla spiaggia di Gili T per affittare i lettini, c'è stata una grande differenza tra chiedere semplicemente "We want four sunbeds" e approcciare con un più amichevole "Hi mate, how are you? What's your name? Nice to meet you Luca, I'm Marco. Can we have four sunbeds, please?" Forse si tratta di una semplice questione di apertura mentale ed educazione, ma è difficile essere aperti e cortesi quando si fatica a trovare le parole per comunicare.

Senza una buona padronanza della lingua non avremmo potuto apprezzare le spiegazioni di Oshi, risolvere gli ordini sbagliati durante l'ultima cena a Gili Asahan, provare il Betel Quid a Sade — quasi certamente Elisa ne sarebbe contenta — o ricevere l'aiuto di Luca/Johan per pulire i lettini dalla sabbia a Gili T. E sono quasi sicuro che a Ubud saremmo finiti in qualche tour truffa.

\begin{center}
    % Decorative line
    \noindent\rule{0.6\textwidth}{0.4pt}
\end{center}

Oltre alla lingua, un'altra competenza fondamentale per un viaggio come questo è la familiarità con la tecnologia. Abbiamo prenotato tutti gli hotel online, gestito i voli con le app di tre diverse compagnie aeree, organizzato gli spostamenti interni con Grab e Gojek, e sincronizzato i biglietti su una cartella nel cloud. Ho pagato quasi sempre tramite le app sul wallet del telefono, comunicato con casa usando una SIM virtuale, e mantenuto attiva la VPN per proteggermi da frodi e furti di dati.

Questa organizzazione ci ha dato una grande libertà di movimento in un viaggio con un programma serrato: siamo riusciti a recuperare un bagaglio perduto senza far slittare le prenotazioni, a tornare sani e salvi a Ubud dal tempio di Besakih, e a organizzare all'ultimo minuto una notte aggiuntiva a Ubud dopo aver scoperto che il vulcano Kawah Ijen era chiuso per attività vulcanica.

Appena tornati a Torino, la prima cosa che ho fatto è stata cancellare dieci applicazioni dal telefono, mantenendone alcune per futuri viaggi. La tecnologia dà un senso di autonomia e indipendenza, ma allo stesso tempo preoccupa pensare a cosa sarebbe potuto accadere se avessi perso il telefono all'inizio della vacanza. Forse è meglio trarre un insegnamento da questa esperienza e, per il prossimo viaggio, dotarsi di un secondo telefono ugualmente configurato come backup.

\begin{center}
    % Decorative line
    \noindent\rule{0.6\textwidth}{0.4pt}
\end{center}

Più scrivo, più mi vengono in mente nuove riflessioni e domande, molte delle quali probabilmente rimarranno senza risposta. Tuttavia, in tutto questo non c'è traccia di malinconia, se non quella legata al fatto di essere tornati a casa. Ciò che abbiamo visto e le persone che abbiamo incontrato rimarranno ricordi indelebili, e spero che il tentativo di trascrivere il più possibile le nostre esperienze aiuti a non dimenticarle. E se questo diario non dovesse bastare, abbiamo quasi diecimila fotografie a supporto, pronte a fornirci ispirazione per rivivere quei momenti.
\newline 


Terima kasih


